% !TITLE 燕山亭·北行见杏花
% !AUTHOR 赵佶
% !DYNASTY 两宋
% !GENRE 词
% !ORDER 11
% !EDITION 《全宋词》，唐圭璋编，中华书局
% !EDITION_STATUS pending

\begin{poem}
裁剪冰绡\textsuperscript{1}，轻叠数重，淡著胭脂匀注。
新样靓妆\textsuperscript{2}，艳溢香融，羞杀蕊珠宫\textsuperscript{3}女。
易得凋零，更多少、无情风雨。
愁苦。问院落凄凉，几番春暮。

凭寄\textsuperscript{4}离恨重重，这双燕，何曾会\textsuperscript{5}人言语。
天遥地远，万水千山，知他故宫\textsuperscript{6}何处。
怎不思量，除梦里、有时曾去。
无据\textsuperscript{7}。和梦也、有时不做。
\end{poem}

\begin{annotations}
\notetext{1}{冰绡：像冰一样洁白透明的薄纱。绡（xiāo），薄的生丝织品。形容杏花花瓣像裁剪好的白色丝绢。}
\notetext{2}{靓妆：美丽的妆饰。靓（jìng），以脂粉妆饰。杏花像化了新式妆容的美人。}
\notetext{3}{蕊珠宫：道教传说中上清的宫阙，仙女居住的地方。杏花的美让蕊珠宫的仙女都自愧不如，极言其美。}
\notetext{4}{凭寄：依靠（燕子）来寄托。想让燕子把离恨带回去，但燕子不懂人言。}
\notetext{5}{会：理解、懂得。燕子怎么会理解人的言语呢——把心事寄托在燕子身上，只是徒劳。}
\notetext{6}{故宫：从前的宫殿，指汴京（今开封）的皇宫。靖康之变后赵佶被金人俘虏北上，再也回不去故宫了。}
\notetext{7}{无据：没有凭据、不可靠。连梦也靠不住——有时候连梦都做不成。连梦中回到故国的机会都被剥夺了，这是最深的绝望。}
\end{annotations}

\begin{appreciation}
他登基时不过十八九岁，写这首词时已是被押往北方的囚徒——词作者赵佶，就是宋徽宗，北宋最后一位皇帝。靖康之变，金兵攻破汴京，徽钦二帝连同后妃宗室一起被掳北上，这首《燕山亭》相传就作于北行途中，见路边杏花而生感，是宋词中写亡国之痛最沉痛的名篇之一。

上阕全在写杏花，却处处是身世。花瓣像裁剪好的洁白冰绡，一层一层轻轻叠起，又像淡淡匀好了胭脂；这样一副新样靓妆，艳光四溢，香气融融，连蕊珠宫的仙女见了都要自愧不如。写到最美处，笔锋骤转：“易得凋零，更多少、无情风雨。”再好的花也架不住风雨，凋谢原是弹指间的事。随后“愁苦。问院落凄凉，几番春暮”——一个被废黜的皇帝，这样的愁苦问的哪里是花，分明是自己：这样的春天，还能有几回。他写花，其实句句都在写自己由盛转衰的那一夜。

下阕写离恨，写得愈发无依。“凭寄离恨重重，这双燕，何曾会人言语。”想托燕子把满腔心事捎回故国，可燕子哪里听得懂人话。“天遥地远，万水千山，知他故宫何处。”山河阻隔，连汴京在哪个方向都辨不清了。“怎不思量，除梦里、有时曾去。”想家的时候，只有做梦能回去一趟。末句更狠：“无据。和梦也、有时不做。”连梦也是靠不住的，有时候连梦都做不成。回家的路先在现实里断掉再在梦里断掉，一层比一层绝望。

王国维在《人间词话》里评此词，说它与李煜同是“以血书者”，但李煜的词如释迦、基督替众生担荷罪孽，赵佶写的却只是一己的身世之痛，境界的大小并不相同。这话说得冷峻，却也公允：一个亡国之君的悔与痛，未必能替他治国时的昏聩开脱。《宋史》给他的评语只有八个字——“诸事皆能，独不能为君耳”。他的字画、他的瘦金体，都是一流；偏偏做了皇帝，而且是一个把江山断送掉的皇帝。

几十年前，王安石还站在金陵城上写下“至今商女，时时犹唱，后庭遗曲”，那是一个清醒者最后的警告，没有人听。几代人过去，警告成了谶语，作词的人成了词里的囚徒。
\end{appreciation}
